\begin{claim}{C3.4}{Proved}{Non-K\"ahler degrees and the precise degreewise statement}
For $k=0,\ldots,6$ the non-K\"ahler degrees are
\begin{equation}\label{eq:C-delta}
(\Delta^k)=(0,2,2+2h,4h,2+2h,2,0).
\end{equation}
If $h\ge1$, the two maps
$\bigoplus_{p+q=k}H^{p,q}_{BC}\to H^k_{\mathrm{dR}}(X;\C)
\to\bigoplus_{p+q=k}H^{p,q}_A$
are both isomorphisms exactly for $k=0,6$.
For the weaker assertion that a $d$-exact pure $d$-closed form of
degree $k$ is $\partial\bar\partial$-exact, the valid degrees are
$k=0,1,6$, and also $k=3,4$ if $h=0$.
\end{claim}
\begin{proof}
The totals from C2.1 are
\[
(h^k_{BC})=(1,0,2,2h,2h,2,1),\qquad
(h^k_A)=(1,2,2h,2h,2,0,1).
\]
Subtracting twice the Betti numbers gives \eqref{eq:C-delta}.
On the two endpoint dots both comparison maps are isomorphisms;
all intermediate de Rham groups vanish.  The first assertion follows
by inspecting when both outside groups vanish.  The weaker assertion
is exactly injectivity of the first comparison map.  In intermediate
degree this is equivalent to vanishing of the Bott--Chern group alone,
which gives the stated list.  In degree zero there is no nonzero
exact form, and in degree six integration identifies Bott--Chern and
de Rham cohomology.  This distinguishes the two meanings of a
``$\partial\bar\partial$-lemma in a degree''; a vanishing Bott--Chern
group in degree one does not make $\Delta^1$ vanish.
\end{proof}
